% =============================================================================
\chapter{Giới thiệu}
% =============================================================================

\section{Bài toán}

Một siêu thị cỡ vừa kinh doanh từ vài nghìn đến hàng trăm nghìn mặt hàng. Mỗi mặt hàng được định danh
bằng \emph{mã vạch} EAN-8 (8 chữ số) hoặc EAN-13 (13 chữ số), thuộc một \emph{nhóm hàng}
(đồ uống, bánh kẹo, gia vị,~\ldots), có đơn giá bán và một ngưỡng tồn kho tối thiểu.

Hàng không được nhập một lần mà theo nhiều \emph{lô}. Mỗi lần nhập tạo một lô mới với số lượng, giá nhập
và \emph{hạn sử dụng (HSD) riêng}: cùng một loại sữa tươi nhưng lô nhập đầu tháng hết hạn sớm hơn lô nhập
cuối tháng. Vì vậy khi bán (xuất kho), siêu thị phải lấy hàng từ lô nhập sớm nhất trước --
nguyên tắc \emph{nhập trước, xuất trước} (FIFO) -- để hạn chế hàng tồn quá hạn. Nếu lô cũ nhất không đủ,
lấy tiếp các lô sau.

Bên cạnh nhập/xuất, người quản lý cần tra cứu nhanh một mặt hàng theo mã vạch, tìm theo tên, xem cảnh báo
hàng sắp hết và hàng sắp hết hạn, nắm giá trị tồn kho theo nhóm và biết mặt hàng nào bán chạy trong tháng.

Đề tài yêu cầu xây dựng chương trình quản lý kho đáp ứng các nhu cầu trên, với trọng tâm là lựa chọn và
tự cài đặt các cấu trúc dữ liệu phù hợp: \textbf{B-Tree bậc $m$} làm chỉ mục mặt hàng theo mã vạch và
\textbf{hàng đợi FIFO trên danh sách liên kết} cho các lô của từng mặt hàng.

\section{Yêu cầu chức năng}

\begin{table}[htbp]
\centering\small
\caption{Yêu cầu của đề tài và mô-đun đáp ứng}
\label{tab:yeucau}
\begin{tabularx}{\linewidth}{@{}clY l@{}}
\toprule
\textbf{STT} & \textbf{Loại} & \textbf{Yêu cầu} & \textbf{Mô-đun} \\
\midrule
1 & Bắt buộc & B-Tree bậc $m$ khóa mã vạch: tìm, chèn (tách nút), xóa (mượn, gộp nút), duyệt & \code{BTree} \\
2 & Bắt buộc & Hàng đợi FIFO các lô cho từng mặt hàng; phiếu nhập tạo lô mới; phiếu xuất lấy từ lô cũ nhất, thiếu thì lấy tiếp lô sau; cập nhật tồn kho; lưu danh sách phiếu & \code{LotQueue}, \code{Phieu}, \code{Inventory} \\
3 & Bắt buộc & Tìm theo tên (chuỗi con), liệt kê theo nhóm; tự cài Merge Sort & \code{Sort}, \code{Report} \\
4 & Bắt buộc & Cảnh báo tồn dưới ngưỡng; cảnh báo hàng hết hạn trong 7 ngày, sắp theo HSD & \code{Report} \\
5 & Bắt buộc & Báo cáo giá trị tồn theo nhóm; top 10 bán chạy trong tháng (theo phiếu xuất) & \code{Report} \\
6 & Bắt buộc & Đọc tệp khi khởi động, ghi khi thoát; kiểm tra dữ liệu nhập (mã trùng, sai kiểu, sai miền) & \code{FileIO}, \code{Validate} \\
7 & Bắt buộc & Dữ liệu mẫu $\ge 30$ bản ghi và bộ lớn $\ge 10^4$; đo thời gian chạy; không rò rỉ bộ nhớ & \code{DataGen}, \code{Benchmark} \\
8 & Nâng cao & In cây sau mỗi lần tách/gộp nút & \code{BTree} \\
9 & Nâng cao & Đếm số nút truy xuất khi tìm với $m = 3, 5, 10$ trên $10^5$ mặt hàng & \code{BTree}, \code{Benchmark} \\
\bottomrule
\end{tabularx}
\end{table}

\paragraph{Ràng buộc cài đặt.} Ngôn ngữ C++17; \textbf{không dùng} \code{std::vector}, \code{std::map},
\code{std::sort} của STL -- mọi cấu trúc (B-Tree, danh sách liên kết, mảng động) và thuật toán sắp xếp đều
tự cài đặt. Dữ liệu lưu trong ba tệp văn bản: \code{mathang.txt}, \code{lo.txt}, \code{phieu.txt}.

\section{Phân công công việc}

Nhóm gồm bốn thành viên. Phần mã nguồn được chia theo mô-đun để làm song song; các kiểu dữ liệu dùng chung
được cả nhóm chốt trong \code{Types.h} ngay từ ngày đầu, mọi thay đổi struct phải báo cả nhóm.

\begin{table}[htbp]
\centering\small
\caption{Phân công mã nguồn}
\label{tab:phancong}
\begin{tabularx}{\linewidth}{@{}l Y l@{}}
\toprule
\textbf{Thành viên} & \textbf{Phụ trách} & \textbf{Tệp nguồn} \\
\midrule
\TVa & B-Tree bậc $m$: tìm, chèn (tách nút), xóa (mượn, gộp nút), duyệt. Nâng cao: in cây sau mỗi lần tách/gộp, đếm số nút truy xuất với $m = 3, 5, 10$ & \code{BTree}, \code{Benchmark} \\
\TVb & Hàng đợi FIFO các lô (DSLK), phiếu nhập/xuất, nhập kho, xuất kho FIFO, cập nhật tồn kho, danh sách phiếu (DSLK) & \code{LotQueue}, \code{Phieu}, \code{Inventory} \\
\TVc & Tìm theo tên, liệt kê theo nhóm, Merge Sort, cảnh báo tồn và HSD, báo cáo giá trị nhóm, top 10 bán chạy & \code{Sort}, \code{Report} \\
\TVd & Kiểu \code{Date}, đọc/ghi tệp, sinh dữ liệu, kiểm tra dữ liệu nhập, menu, đo thời gian, kiểm tra rò rỉ bộ nhớ & \code{Date}, \code{FileIO}, \code{DataGen}, \code{Validate}, \code{main}, \code{Benchmark} \\
\bottomrule
\end{tabularx}
\end{table}

\begin{table}[htbp]
\centering\small
\caption{Phân công báo cáo và trình bày}
\label{tab:phancongbaocao}
\begin{tabular}{@{}ll@{}}
\toprule
\textbf{Phần} & \textbf{Người thực hiện} \\
\midrule
Chương 1. Giới thiệu, bảng phân công & \TVd \\
Chương 2. Phân tích thiết kế & Cả nhóm, \TVa{} tổng hợp \\
Chương 3. Giải thuật & Mỗi người viết phần mình cài đặt \\
Chương 4. Cài đặt và kiểm thử & \TVd{} (mỗi người gửi test case phần mình) \\
Chương 5. Thực nghiệm & \TVa{} và \TVd \\
Chương 6. Kết luận, tài liệu tham khảo & \TVc \\
Slide thuyết trình & \TVc{} làm chính, mỗi người 2--3 slide \\
\bottomrule
\end{tabular}
\end{table}

\paragraph{Kế hoạch.} Ngày 1 chốt \code{Types.h}, tên hàm và định dạng tệp; ngày 2--4 cài đặt song song
(B-Tree: tìm và chèn ngày 2--3, xóa ngày 4); ngày 5 ghép vào menu, sửa lỗi, kiểm tra rò rỉ bộ nhớ;
ngày 6 chạy thực nghiệm, phần nâng cao, hoàn thiện báo cáo; ngày 7 đọc chéo mã nguồn và tập vấn đáp.
Mỗi tối các thành viên đẩy mã lên kho Git chung và ghép thử.

\section{Cấu trúc báo cáo}

Chương~\ref{chap:thietke} phân tích thiết kế: mô hình dữ liệu, lý do chọn từng cấu trúc so với phương án khác,
sơ đồ mô-đun. Chương~\ref{chap:giaithuat} trình bày giải thuật với mã giả, ví dụ chạy tay và độ phức tạp.
Chương~\ref{chap:caidat} mô tả cài đặt, giao diện và kiểm thử. Chương~\ref{chap:thucnghiem} trình bày thực
nghiệm so sánh B-Tree với $m = 3, 5, 10$. Chương~\ref{chap:ketluan} kết luận và nêu hướng phát triển.
